% !TeX root = ../main.tex
\section{Выводы}

В ходе выполнения практической работы было проведено имитационное моделирование сквозного цифрового тракта пространственно-временной обработки радиолокационных сигналов.%[cite: 6] На основе полученных результатов можно сделать следующие выводы:

\begin{enumerate}
    \item \textbf{Работоспособность сквозного тракта.} Успешно реализована и отлажена цепочка преобразований, включающая аналого-цифровое преобразование, цифровой перенос частоты (ЦДМ), низкочастотную фильтрацию побочных составляющих (ЦФ), децимацию и согласованную фильтрацию (СФ).%[cite: 6] Разработанный модульный алгоритм корректно извлекает полезный сигнал из шума и производит его оптимальное сжатие по времени.%[cite: 6]
    
    \item \textbf{Субдискретизация и перенос частоты.} Использование частоты дискретизации АЦП ($140$ МГц), которая меньше промежуточной частоты сигнала ($170$ МГц), на практике подтвердило эффективность метода субдискретизации (bandpass sampling).%[cite: 6] Возникающий при этом эффект алиасинга носит управляемый характер: цифровой гетеродин успешно компенсирует разностную частоту и переносит спектр сигнала точно в нулевую (baseband) область частот.%[cite: 6]
    
    \item \textbf{Вычислительная эффективность.} Применение метода сегментированной быстрой свертки на основе алгоритма перекрытия со сложением (реализованного функцией \texttt{fftfilt}) позволяет математически точно и эффективно фильтровать длинные последовательности данных.%[cite: 4, 6] В отличие от вычисления линейной свертки во временной области, блочная обработка через прямое и обратное преобразования Фурье исключает переполнение памяти ЭВМ и радикально снижает требуемое количество вычислительных операций.%[cite: 1, 4]
    
    \item \textbf{Влияние оконного взвешивания.} Сравнительный анализ откликов согласованного фильтра показал, что применение весовой функции (окна Хэмминга) к опорному сигналу позволяет существенно подавить уровень боковых лепестков сжатого радиоимпульса.%[cite: 6, 7, 15] Несмотря на неизбежное побочное действие в виде незначительного расширения главного лепестка, применение окна является необходимым в радиолокации для предотвращения маскировки слабых эхо-сигналов боковыми лепестками от мощных целей.%[cite: 15]
\end{enumerate}